\documentclass[11pt,a4paper]{amsart}
\usepackage[a4paper,margin=30mm]{geometry}
\usepackage[T1]{fontenc}
\usepackage{lmodern,amssymb,mathtools,microtype}
\usepackage[hidelinks]{hyperref}
\hypersetup{pdftitle={Discrepancy bounds for Kritzinger's greedy sequence},pdfauthor={Lazar Korsić}}
\newtheorem{theorem}{Theorem}[section]
\newtheorem{lemma}[theorem]{Lemma}
\newtheorem{corollary}[theorem]{Corollary}
\newcommand{\norm}[1]{\lVert #1\rVert}
\newcommand{\ind}{\mathbf{1}}
\title{Discrepancy bounds for Kritzinger's greedy sequence}
\author{Lazar Korsi\'c}
\date{November 23, 2026}
\subjclass[2020]{11K38}
\keywords{Star discrepancy, greedy sequences, integrated discrepancy}
\begin{document}
\begin{abstract}
Kritzinger constructed a sequence in the unit interval by choosing each new
point to minimize the $L^2$ discrepancy of the enlarged set. We prove that
its unnormalized star discrepancy is at most $\log(2N)+3/2$ for every
prefix of length $N$, establishing the order assertion in his second
conjecture. The proof traces averages of the discrepancy backwards through
the insertions that change one endpoint of a test interval. A second argument
bounds the height of a chord above the primitive of the discrepancy. It
gives a uniform bound on that primitive and an $O_p(\sqrt{\log(2N)})$
bound on the $L^p$ discrepancy for every fixed $1\le p<\infty$.
The estimates hold for every choice among global minimizers and extend,
with modified constants, to arbitrary finite initial configurations.
\end{abstract}
\maketitle

\section{Introduction}
For points $x_1,\ldots,x_N$ in $[0,1]$ let
\[
 D_N(t)=\#\{i\le N:x_i<t\}-Nt,\qquad H_N(t)=\int_0^t D_N(u)\,du .
\]
We use unnormalized discrepancy throughout. Thus a value of $D_N(t)$
measures the error in the \emph{number} of points in $[0,t)$.
Kritzinger's Algorithm~2 \cite[\S4.1]{K} starts at $x_1=1/2$ and
chooses the leftmost minimizer of $\int_0^1 D_{N+1}(t)^2\,dt$ as the next
point. Equivalently, it minimizes
\begin{equation}\label{eq:rule}
 (N+1)x^2-x-2\sum_{i=1}^N\max(x,x_i),\qquad 0\le x\le1 .
\end{equation}
We call any choice of a global minimizer a greedy insertion. Starting
from the empty set gives $x_1=1/2$ by the same rule. Steinerberger
\cite{S} identified this construction with greedy minimization of the
quadratic Wasserstein distance to uniform measure. He also proved that
$\norm{H_N}_\infty\le 2N^{1/3}$ for infinitely many $N$ and asked
whether a constant bound could hold.

\begin{theorem}\label{th:main}
For every greedy sequence from the empty set and every $N\ge1$,
\[
 \norm{D_N}_\infty\le \log(2N)+\tfrac32,\qquad
 \norm{H_N}_\infty\le (1+\sqrt5)/4 .
\]
For every fixed $1\le p<\infty$ one also has
$\norm{D_N}_{L^p[0,1]}\ll_p\sqrt{\log(2N)}$.
\end{theorem}
Here $\log$ is natural. The first estimate proves the $O(\log N)$
assertion in Kritzinger's Conjecture~2. His further prediction that
$\limsup_{N\to\infty}\norm{D_N}_\infty/\log N<1/(3\log 2)$
remains open. The case $p=2$ gives the growth order sought in
Conjecture~1; its strict comparison with the symmetrized van der Corput
sequence also remains open \cite[\S4.1]{K}.

At each insertion the greedy rule bounds certain interval averages
on both sides of the new point. The difficulty is that these bounds
apply at different points at different times. Section~\ref{sec:trace}
fixes a threshold and follows it backwards, shortening its test interval
whenever it encounters an interior insertion. The accumulated cost is
logarithmic in the ratio of the initial to the final length.
Sections~\ref{sec:chord} and~\ref{sec:mom} instead retain both endpoints
and control the interpolation error of $H_N$. This gives an estimate
on the frequency of large differences between neighboring cell averages,
which suffices for the finite moments.

\paragraph{Provenance and author responsibility.}
I asked OpenAI's ChatGPT models to investigate Kritzinger's problem.
The models carried out the mathematical investigation, developed the
proofs, wrote and ran the computational checks, and drafted this
manuscript. My role is to review the arguments and references and to
take responsibility for the version I release. This draft precedes
completion of my review.
% After personal verification, the author may replace the last two sentences by:
% I have checked the arguments and references and take responsibility for this version.

\section{What global minimization implies}\label{sec:cone}
Write $S_m(a,b)=(H_m(b)-H_m(a))/(b-a)$ for the mean discrepancy on
$[a,b]$, and set $K_m(t)=H_m(t)+t(1-t)/2$. The identity
\[
 H_m(t)=\sum_{i\le m}(t-x_i)_+-mt^2/2
\]
shows that \eqref{eq:rule} equals $-2K_N(x)-2\sum_{i\le N}x_i$.
Thus $c=x_m$ maximizes $K_{m-1}$.

\begin{lemma}\label{lem:cone}
At the insertion of $c=x_m$ one has
$H_m(u)-H_m(c)\le |u-c|/2$ for every $u$. In particular,
\begin{equation}\label{eq:slopes}
 S_m(a,c)\ge-1/2,\qquad S_m(c,b)\le1/2\qquad(a<c<b).
\end{equation}
For every canonical-start prefix, $H_m(t)\le t/2$ and
$H_m(t)-H_m(1)\le(1-t)/2$.
\end{lemma}
\begin{proof}
Adding $(u-c)_+-u^2/2$ to $H_{m-1}(u)$ gives
\[
 H_m(u)-H_m(c)=K_{m-1}(u)-K_{m-1}(c)+|u-c|/2 .
\]
Maximality proves the first assertions. For the boundary estimate induct
from $H_0=0$. At $t\le c$ the update is nonpositive. At $t>c$ the cone
gives $H_m(t)\le H_m(c)+(t-c)/2\le t/2$. Under reflection $y_i=1-x_i$
the new primitive is $H_m(1-u)-H_m(1)$, by integration. Since $u(1-u)$
is symmetric, reflection preserves global maximizers and gives the other
boundary estimate.
\end{proof}

We recall the rank and separation properties of \cite[Theorems~2--3]{K}
in a form that includes ties. Sort the $m-1$ old points as $y_i$ and
put $A_j=\sum_{i=1}^j y_i$, $A_0=0$. The hinge sum is
$\max_{0\le j<m}(ju-A_j)$. The corresponding quadratics in $K_{m-1}$
have vertices $(j+1/2)/m$ and vertex heights
$P_j+1/(8m)$, where $P_j=j(j+1)/(2m)-A_j$.
All vertices lie in $(0,1)$, so every global maximizer is such a vertex
with largest $P_j$. Comparison with its neighboring indices yields
$y_j\le j/m$ and $y_{j+1}\ge(j+1)/m$ whenever the corresponding
neighbor exists. Every new point is therefore at distance at least
$1/(2m)$ from every old point. In particular,
\begin{equation}\label{eq:spacing}
 |x_i-x_j|\ge 1/(2N)\qquad(1\le i<j\le N)
\end{equation}
for the empty start.

\section{Tracing an interval backwards}\label{sec:trace}
Fix $0\le a<t\le1$. The quantity to control is the excess of the
discrepancy at the threshold over its average immediately to the left:
\[
 R_m(a,t)=D_m(t)-S_m(a,t)
   =\sum_{\substack{i\le m\\a<x_i<t}}\frac{x_i-a}{t-a}
       -\frac{m(t-a)}2 .
\]
The second equality follows by integrating each counting step. Both
endpoints are excluded from the sum, including when $t$ is an old point.

\begin{lemma}\label{lem:trace}
For the empty start,
$R_N(a,t)\le 1+\max(0,\log(2N(t-a)))$.
\end{lemma}
\begin{proof}
Let $h=t-a$. If $c=x_m$ is outside $(a,t)$, the displayed formula gives
$R_m(a,t)=R_{m-1}(a,t)-h/2$. If $a<c<t$, put $h'=t-c$. Splitting the
average at $c$ and applying \eqref{eq:slopes} gives
\[
 S_m(c,t)-S_m(a,t)=\frac{c-a}{h}\big[S_m(c,t)-S_m(a,c)\big]
 \le 1-h'/h .
\]
Subtract the two means from the same value $D_m(t)$. On the shortened
interval $c$ is now its left endpoint, hence
\[
 R_m(a,t)\le R_{m-1}(c,t)-h'/2+1-h'/h .
\]
Iterate towards time zero. Discarding the nonpositive terms, we obtain
\[
 R_N(a,t)\le \sum_{i=1}^k(1-h_i/h_{i-1}),
\]
where $h_0=t-a$ and $h_i$ is the length after the $i$th shortening.
There is no terminal contribution since $R_0=0$.

For $k\le1$ the sum is at most one. Otherwise the last two accepted
points are $t-h_{k-1}$ and $t-h_k$. By \eqref{eq:spacing},
$h_{k-1}\ge h_{k-1}-h_k\ge1/(2N)$. Charge the last shortening at
most one and sum $1-v\le-\log v$ over the others. Their sum is at most
$\log(h_0/h_{k-1})\le\log(2Nh_0)$. Isolating the last term allows
$t$ to lie arbitrarily close to a point.
\end{proof}
Take $a=0$. Lemma~\ref{lem:cone} bounds the initial-interval mean by
$1/2$, which proves the upper estimate on $D_N(t)$ in
Theorem~\ref{th:main}. Off the support, the discrepancy of the reflected
configuration at $1-t$ is $-D_N(t)$. This gives the lower estimate.
Limits from the complement of the support give both one-sided limits
at each jump, and $D_N(0)=D_N(1)=0$ for the empty start.

\section{A bound on chord heights}\label{sec:chord}
For $a<t<b$ set $p=t-a$, $q=b-t$ and write
\[
 G_m(a,t,b)=\frac{qH_m(a)+pH_m(b)}{p+q}-H_m(t).
\]
\begin{theorem}\label{th:chord}
For all such triples and all greedy prefixes from the empty set,
$G_m(a,t,b)\le\sqrt{pq}$.
\end{theorem}
\begin{proof}
Induct from $G_0=0$. When $c=x_m$ lies outside $(a,b)$, the new hinge
is affine on $[a,b]$. The quadratic update gives
$G_m=G_{m-1}-pq/2$ there. For an interior insertion, applying the cone
at both endpoints yields
\begin{equation}\label{eq:chordbirth}
 G_m(a,c,b)\le (c-a)(b-c)/(b-a).
\end{equation}
This suffices if $c=t$. If $a<c<t$, interpolate first at $c$ to get
\[
 G_m(a,t,b)=G_m(c,t,b)+\frac{q}{b-c}G_m(a,c,b).
\]
Put $p'=t-c$. The insertion is an endpoint of the first interval, so
the first term is at most $\sqrt{p'q}-p'q/2$. Use
\eqref{eq:chordbirth} for the second term. It remains to check
\[
 \sqrt{p'q}+\frac{q(p-p')}{p+q}\le\sqrt{pq}.
\]
Indeed $\sqrt q/(\sqrt p+\sqrt{p'})\ge \sqrt q/(2\sqrt p)
\ge q/(p+q)$. Multiplying by $p-p'$ proves the inequality.
The case $t<c<b$ follows by reflection.
\end{proof}

For $I=[a,b]$ define the tent
$\psi_I(x)=\max(0,\min(x-a,b-x))/(b-a)$. At its midpoint $t$,
the hinge formula gives
\[
 Q_N(I):=\sum_{i\le N}\psi_I(x_i)-N|I|/4
       =2G_N(a,t,b)/|I|\le1 .
\]
At $a=0$, $b=1$, Theorem~\ref{th:chord} gives
$H_N(t)\ge tH_N(1)-\sqrt{t(1-t)}$.
Lemma~\ref{lem:cone} gives $|H_N(1)|\le1/2$ and $H_N(t)\le t/2$.
Writing $u=2t-1$, Cauchy--Schwarz gives
\[
 t/2+\sqrt{t(1-t)}=(1+u+2\sqrt{1-u^2})/4\le(1+\sqrt5)/4 .
\]
This proves the primitive estimate in Theorem~\ref{th:main}.

\section{The finite moments}\label{sec:mom}
Partition $[0,1]$ into $M$ equal cells. Write $s_i$ for the mean of
$D_N$ on cell $i$, and $z_i=s_{i+1}-s_i$ for $1\le i<M$.
The midpoint estimate gives $z_i\le2$.
A negative jump has no analogous pointwise bound from
Theorem~\ref{th:chord}. The following estimate on its frequency is sufficient.
\begin{lemma}\label{lem:tail}
For $A>0$, $\#\{i:z_i<-A\}/M\le\min(1,2^{(5-A)/4})$.
\end{lemma}
\begin{proof}
For any interval split into lengths $p,q$, its mean minus its left
mean is $G_N/p$, and its right mean minus its mean is $G_N/q$.
Thus extending a cell to the right by at most its length increases its
mean by at most one. Extending it to the left by at most its length
decreases its mean by at most one. For a run of $m$ cells, grow its
first cell by successive doublings towards the right, stopping at the
end of the run. Repeat from its last cell towards the left. This proves
\[
 s_{\rm last}-s_{\rm first}\le 2\lceil\log_2 m\rceil .
\]
Cut the row at the $k$ jumps below $-A$. If $k>0$, the $k+1$ resulting
run lengths sum to $M$. The boundary estimates give $s_M-s_1\ge-1$.
Summing the rises along runs and the drops across cuts, then using
concavity of the logarithm, yields
\[
 -1\le-kA+2(k+1)\big[1+\log_2(M/(k+1))\big].
\]
Divide by $k$ and use $(k+1)/k\le2$. This gives
$A\le5+4\log_2(M/(k+1))$, as required. The case $k=0$ is immediate.
\end{proof}
Integrating the tail against $pA^{p-1}\,dA$ with $p=d\ge1$, and using $z_i\le2$,
shows that $M^{-1}\sum_{i<M}|z_i|^d\le C_d$ uniformly in $M,N$.

Let $U_j$ be the function of dyadic cell means at level $j$. On a
parent cell the two child means are $U_j-V_j$ and $U_j+V_j$, where
$V_j$ is half the right-minus-left jump. In the grid of $M=2^{j+1}$
cells these are precisely the jumps of odd index. Hence
\[
 \int_0^1 |V_j|^d=\frac{2}{M}\sum_{i\ {\rm odd},\,i<M}|z_i/2|^d
 \le 2^{1-d} C_d .
\]
For $d=2$ the two signs cancel the cross term, and each level adds
at most a constant to $\norm{U_j}_2^2$. For an even integer $d=2r$,
put $u=\norm{U_j}_d$ and $v=\sup_j\norm{V_j}_d$. Average the
binomial expansions on the two children and apply H\"older's inequality
on the parents. Only even powers of $V_j$ remain, so
\[
 \norm{U_{j+1}}_d^d
 \le\sum_{\ell=0}^r\binom{2r}{2\ell}u^{2r-2\ell}v^{2\ell}
 \le (u^2+2^{2r}v^2)^r .
\]
The last bound follows coefficient by coefficient: for $\ell\ge1$,
$\binom{2r}{2\ell}\le 2^{2r}\le\binom r\ell 2^{2r\ell}$.
Taking the $r$th root and summing gives
$\norm{U_L}_d^2\le1/4+L2^{2r}v^2$.
Take $L=\lceil\log_2(4N)\rceil$. By \eqref{eq:spacing} the interior
of each final cell contains at most one point, so the essential
oscillation there is at most $1+N2^{-L}\le5/4$.
Consequently $\norm{D_N-U_L}_\infty\le5/4$ in essential supremum.
For any fixed $1\le p<\infty$, take an even $d\ge p$ and use
$\norm{D_N}_p\le\norm{D_N}_d$. This completes Theorem~\ref{th:main}.

\section{Initial configurations and quadrature}
\begin{corollary}
Start with any list of $s$ points in $[0,1]$, allowing repeats and
endpoints, and extend it by global greedy insertions. Put
$B=\max(1/2,\norm{D_s}_\infty)$ and $C=\max(1,B)$. For $N\ge\max(1,s)$,
\[
 \norm{D_N}_\infty\le3B+1+\log(2N),\quad
 \norm{H_N}_\infty\le B+C/2,
\]
and, for every fixed $1\le p<\infty$,
\[
 \norm{D_N}_p\ll_{s,p}\sqrt{\log(2N)}.
\]
\end{corollary}
\begin{proof}
At time $s$ the boundary means in either orientation are at most $B$
at the left endpoint. The induction of Lemma~\ref{lem:cone} propagates
this since $B\ge1/2$. Stop the backward trace at $s$, where the final
excess is at most $2B$. All accepted points in the trace are post-seed
insertions, for which the spacing argument still applies. Adding the
boundary mean gives the star bound. Reflection holds almost everywhere;
inward limits recover the required endpoint convention even with atoms
at $0$ or $1$.

At time $s$, $G_s$ equals $pq/(p+q)$ times the difference of the right
and left means, so $G_s\le 2Bpq/(p+q)\le C\sqrt{pq}$. The induction
of Theorem~\ref{th:chord} preserves this bound since $C\ge1$.
On $[0,1]$ the chord estimate and $|H_N(1)|\le B$ give the primitive
bound. Apply the proof of Lemma~\ref{lem:tail} to means divided by $C$:
the rise estimate is unchanged and the endpoint difference is at least
$-2B/C\ge-2$. This replaces 5 by 6 in its exponent numerator.
The moment argument carries through, with final-cell oscillation at most
$s+1+1/4$ because only the post-seed points require separation.
\end{proof}

For the empty start the cone also has a quadrature consequence. Write
$E(f)=\sum_{i\le N} f(x_i)-N\int_0^1 f$. Then $|E(x)|=|H_N(1)|\le1/2$
for the identity function and
$E(|x-t|)=2H_N(t)-H_N(1)\le1/2$ by the two boundary bounds.
Every convex $L$-Lipschitz function has the representation
$f(x)=b+ax+\int_{(0,1)}|x-t|\,d\mu(t)$, where $\mu=f''/2\ge0$
in the distributional sense and $a$ is the average of the two endpoint
slopes. Since $|a|+\mu((0,1))\le L$, it follows that $E(f)\le L/2$.

The same primitive gives an exact balance for the squared discrepancy.
Let $\mathcal E_m=\int_0^1D_m^2$, $M_m=\max K_m$ and
$A_m=\int_0^1K_m$. Then
\[
 \mathcal E_{m+1}-\mathcal E_m=\tfrac16-2(M_m-A_m).
\]
Indeed the increment in $D_m(u)$ is $\ind_{\{u>c\}}-u$ at
$c=x_{m+1}$. Its squared integral is $c^2-c+1/3$, and integration by
parts gives $\int_0^1D_m(u)(\ind_{\{u>c\}}-u)\,du=\int_0^1 H_m-H_m(c)$.
Substitution of $K_m(c)=M_m$ gives the formula.

\section{A sharper tent bound}
The conjectured inequality $-1\le Q_N(I)\le1/2$ on dyadic intervals
would sharpen the one-sided estimate of Section~\ref{sec:chord}.
There is a finite test for an even stronger upper bound. Extend the tent
to arbitrary center $z\in\mathbb R$ and radius $r>0$ by setting
\[
 T(z,r)=\tfrac12\sum_{i\le N}(1-|x_i-z|/r)_+-Nr/2 .
\]
Then $T(z,r)\le1/2$ for all $z,r$ if and only if every consecutive
block $J$ of $k\ge1$ sorted points satisfies
\[
 4N\sum_{x\in J}|x-\operatorname{med}(J)|\ge(k-1)^2 .
\]
Indeed, twice the tent value is the maximum over subsets $J$ of
$|J|-\sum_{x\in J}|x-z|/r-Nr$. For fixed nonempty $J$ choose a median
to minimize the deviation sum $b_J$; the supremum over radii is then
$|J|-2\sqrt{Nb_J}$, including $b_J=0$ by a limit. The empty subset
has supremum zero. At each size $k$ the least deviation is attained on
a consecutive block: take the $k$ points nearest an optimizing median
and then optimize their median. This proves the equivalence.

Exact integer comparisons in the accompanying code verified the block
inequality for every canonical prefix through $N=512$ and every global
tie history through $N=24$ (1,798 distinct terminal configurations).
The proof of Theorem~\ref{th:main} is independent of these finite checks.

\begin{thebibliography}{9}
\bibitem{K} R.~Kritzinger, Uniformly distributed sequences generated by a
greedy minimization of the $L_2$ discrepancy, \emph{Moscow J. Combin.
Number Theory} \textbf{11} (2022), 215--236.
\href{https://arxiv.org/abs/2109.06298}{arXiv:2109.06298}.
\bibitem{S} S.~Steinerberger, On combinatorial properties of greedy
Wasserstein minimization, \emph{J. Math. Anal. Appl.} \textbf{532}
(2024), article 127940.
\href{https://arxiv.org/abs/2207.08043}{arXiv:2207.08043}.
\end{thebibliography}
\end{document}
